\section{Recognizing and reconstructing deployed polygons}
\label{sec:polygons}
This section determines exactly which signed boundary words arise from
an additional pseudoline, and recovers the weighted incidences from
those words. Local triangle tests alone do not determine the image.

Fix the labelled old arrangement, orient its lines, and write
$E=\{1,\ldots,n\}$. Its chirotope $\chi_M(i,j,k)$ records the sign
of an ordered triple of old lines, with zero for a concurrent triple.
For normal vectors $a_i$, our convention is
$\chi_M(i,j,k)=\operatorname{sgn}\det(a_i,a_j,a_k)$; a line with
affine normal $(u_i,v_i,w_i)$ has positive endpoint direction
$(-v_i,u_i)$. Let $d$ denote the additional line and put
$q(i,j)=\chi_N(i,j,d)$ for the extended chirotope.

We use the rank-three topological representation theorem in its
projective form~\cite[Section~7.1.4, pp.~6--7]{RGZ}.
Equal labelled covector data give equal cell decompositions on the
spherical double covers. They can be matched equivariantly: match
vertices, extend over each pair of antipodal edges, and then extend
the boundary homeomorphism over each pair of two-dimensional cells.
These cells are disks. This also follows by cutting along one old
line and adding the remaining proper arcs successively: each arc
segment cuts a disk face into disks. Choosing the map on one member
of each antipodal pair determines it on the other. The resulting
homeomorphism descends to the projective arrangements. Thus a
topological representation of an extension can always be carried to
the fixed old arrangement; its deletion is not allowed to vary.

A signed boundary word is a cyclic list $W$ of
$\widetilde E=\{+i,-i:i\in E\}$, each label occurring once.
Temporarily choose its cyclic orientation. Set $q_W(i,j)=0$ when
the endpoint pairs of $i$ and $j$ do not alternate. Otherwise set
$q_W(i,j)=+1$ if $+j$ lies in the positive open arc from $+i$ to
$-i$, and set it to $-1$ otherwise. Put $q_W(i,i)=0$.
This gives an alternating sign map even for a word that is not
realizable in the prescribed old arrangement.

For any alternating $q$ define
\begin{equation}\label{eq:polygon-combined}
 \chi_q(i,j,k)=\chi_M(i,j,k),\qquad
 \chi_q(i,j,d)=q(i,j),
\end{equation}
and extend by alternation, with zero on repeated entries.
The complete finite chirotope test has two parts
\cite[Section~7.2.3, p.~10]{RGZ}.
First, its nonzero unordered triples must satisfy the matroid basis
exchange axiom: for bases $B,C$ and $b\in B\setminus C$ there is
$c\in C\setminus B$ with $(B\setminus\{b\})\cup\{c\}$ a basis.
Second, for every five distinct labels $\lambda,a,b,c,e$, the three
signs
\begin{equation}\label{eq:polygon-gp}
 \chi_q(\lambda,a,b)\chi_q(\lambda,c,e),\quad
 -\chi_q(\lambda,a,c)\chi_q(\lambda,b,e),\quad
 \chi_q(\lambda,a,e)\chi_q(\lambda,b,c)
\end{equation}
must contain both signs or be all zero. The basis condition is
necessary when zeros occur.

\subsection{Disk systems, injectivity and packet order}
The disk criterion explains why opposite packets appear, including
the degenerate case in which every pair of cut arcs is disjoint.

\begin{proposition}[The proper-arc disk criterion]\label{prop:disk-criterion}

Let $\Delta$ be a closed disk, with finitely many labelled tame proper
arcs $\delta_i$ whose $2n$ endpoints on $\partial\Delta$ are distinct.
Suppose each pair meets at most once, at an interior transverse
crossing, and disjointness together with equality is an equivalence
relation on the arcs.  If there are at least two equivalence classes,
this system is the cut-open system of a projective pseudoline
arrangement and an added pseudoline.  If there is just one class,
the same assertion holds exactly when its endpoints form two
consecutive packets, each containing one endpoint of each chord.
The reconstructed old arrangement is non-pencil exactly when its
pairwise intersection points are not all the same point.
\end{proposition}
\begin{proof}

We first prove that the packet condition is automatic with at least
two classes.  Fix a class $C$ and an arc $\beta$ in another class.
Every arc in $C$ crosses $\beta$ once.  The pairwise disjoint arcs of
$C$ therefore separate the two endpoints of $\beta$ and occur in a
linear chain along it.  To see explicitly that there is no branch,
cut $\Delta$ along $\beta$.  In either resulting half-disk the pieces
of arcs in $C$ are pairwise disjoint arcs joining distinct points of
the cut to the opposite boundary; noncrossing forces the same linear
order, reversed when the boundary is traversed in the opposite
direction.  Thus the endpoints of $C$ occur in two packets with
opposite orders.

No endpoint of an arc outside $C$ can lie strictly between two
consecutive endpoints in either packet.  Indeed those two arcs of
$C$ bound a middle strip.  An arc starting in that strip and crossing
each of its two boundary arcs exactly once cannot end in a component
outside both: after the first crossing it is outside one boundary,
and to reach the other boundary it would have to recross the first.
This contradicts the one-crossing hypothesis.  Since every outside
arc crosses both boundaries, neither endpoint can start in the
middle strip.  It follows that the two packets are consecutive among
all chord endpoints, not merely among endpoints belonging to $C$.

Collapse each packet to a point of the boundary.  The quotient of the
disk is again a disk.  For completeness, one may describe the quotient
map in polar coordinates: choose a nondecreasing degree-one boundary
map $\phi$ that is constant exactly on the packet intervals, and use
$t\mapsto (1-r)t+r\phi(t)$ on the circle of radius $r<1$.  These maps
are strictly increasing for $r<1$, extend continuously to the boundary
collapse, and induce a continuous bijection from the quotient onto a
disk.  Compactness gives a homeomorphism.  Finite piecewise linear
choices give the corresponding tame local model.

Choose a representative chord from each of the $r$ classes.  They
cross pairwise.  Their $2r$ boundary points have order
$C_1,\ldots,C_r,C_1,\ldots,C_r$: after cutting at the two endpoints
of $C_1$, every other chord has one endpoint in either half, and two
such chords cross exactly when their orders in the two halves agree.
This proves the displayed order.  Consequently the involution pairing
the two points of each class is an order-preserving, fixed-point-free
boundary involution.  Identify this boundary by that involution.
After a boundary homeomorphism it is the usual antipodal
identification, so the quotient is $\mathbb {RP}^2$.

The image of the boundary is a new pseudoline $d$.  Each old chord
closes to a simple curve, with its two endpoints identified and no
other self-identification.  It crosses $d$ once, hence is an essential
one-sided simple curve and is a pseudoline.  Chords from different
classes retain their unique interior transverse intersection.  Chords
in the same class meet only at the boundary vertex created for that
class.  Their opposite packet orders give, in the local disk after
gluing, the cyclic order of rays of a pencil: each pair of old curves
crosses there, and each crosses $d$.  Thus all required pairwise
intersection and crossing conditions hold.  Expanding the packet
vertices inside a narrow band about $d$ recovers the original disk
system, up to a labelled homeomorphism.

For one class the same construction works when two packets exist.
Conversely, if every old chord is disjoint, every pair of old
pseudolines must have met on $d$.  Since each old pseudoline meets
$d$ only once, all those intersections are a single vertex, and a
small neighborhood of that vertex gives exactly two packets.  The
condition is necessary.  Finally the quotient construction lists all
old intersection points, interior crossings and class vertices;
their being a single common point is precisely the pencil case.
\end{proof}

\begin{proposition}[A signed polygon determines its position]\label{prop:polygon-injective}

For the fixed labelled non-pencil arrangement $\Sigma$, the map from
ordinary relative positions to literal deployed signed
polygons is injective.  The polygon determines the rank-two
contraction by the signed pair-alternation rule, including all its
parallel classes.  The only ambiguity is reorientation of the added
element, corresponding to reversal of the chosen cyclic orientation.
\end{proposition}
\begin{proof}

The Jordan separation argument used to define $q_W$ recovers every value
$q(i,j)=\chi_N(i,j,d)$ from the word.  In particular, no guess at a
cluster decomposition is needed: $i$ and $j$ are in the same cluster
if and only if their endpoint pairs do not alternate.  For an actual
word this is exactly equality of their intersections with $d$, hence
is an equivalence relation.  We know $\chi_N$ on every old triple
from $M$, and on every triple containing $d$ from $q$.  Alternation
and zeros on repetitions determine all remaining ordered triples.
Thus \eqref{eq:polygon-combined} determines the entire extension.  The
relative topological representation comparison identifies its
relative position in the fixed arrangement.  Reversing $q$ changes
exactly the chirotope entries containing $d$, which is precisely
reorientation of $d$.

The contraction determines the extension because its entries and
$\chi_M$ specify every ordered triple. The packet order contains
additional literal boundary information, determined next.
\end{proof}

For the explicit packet expansion, suppose that $\chi_q$ passes the
chirotope test and that $q$ has no zero row. Its rank-two oriented
matroid has a cyclic order of signed parallel-ray classes. Within a
class containing $\varepsilon i$ and $\eta j$, choose $k$ outside
the corresponding unoriented parallel class and order the two labels by
\begin{equation}\label{eq:polygon-expansion}
 \varepsilon i\text{ precedes }\eta j
 \quad\Longleftrightarrow\quad
 -\eta\chi_M(i,j,k)q(i,k)=+1.
\end{equation}
Place the opposite class in reversed, sign-negated order.
The resulting word is denoted $\operatorname{Exp}_M(q)$.
The next proposition proves that this prescription is independent of
$k$ and gives a total order in every packet.

\begin{proposition}[The exact packet reversal]\label{prop:packet-order}

If the vertices of $\Sigma$ on $d$ have line clusters of sizes
$k_1,\ldots,k_\nu$, the deployed word consists of paired consecutive
packets in opposite order.  Let $H$ be the half-turn on the $2n$
positions and let $R$ reverse the order separately within every
packet, fixing singleton packets.  The label-negation involution is
$H\circ R=R\circ H$.  It equals $H$ exactly when $d$ is stable.
The within-packet order is \eqref{eq:polygon-expansion}.
\end{proposition}
\begin{proof}

At an old vertex $v$ on $d$, choose a small disk containing no other
old vertex.  The $k$ old line germs and the germ of $d$ form a pencil
of transverse proper arcs.  A homeomorphism of this disk carries
the finite star to straight radial segments, preserving its cyclic
order.  Take $d$ locally to be $y=0$ and write the other lines as
$x=b_i y$, with distinct $b_i$.  The two nearby boundary arcs meet
them at $x=\epsilon b_i$ and $x=-\epsilon b_i$.  Their orders are
reversed.  This proof is topological: the straightening is only of
one local star, and is available whether the whole arrangement is
stretchable or not.

After the packets are collapsed, the endpoints are the rank-two
cyclic order on $d$.  Corresponding packets are opposite and have
the same cardinality.  Hence the full word has the form
\[
       B_1\cdots B_r\,(-B_1)^{\rm rev}\cdots(-B_r)^{\rm rev}.
\]
Its first and second halves each have $n$ entries.  The formula for
label negation follows entry by entry.  If a packet has at least two
entries, its reversal is nonidentity, so label negation is not the
half-turn.  If every packet is a singleton, the reversals are all
identities.  Thus the equivalence with stability is exact.

We also verify the signed expansion formula, needed for the
algorithm.  It involves only two lines $i,j$ of one cluster, $d$,
and an old line $k$ outside that cluster.  This pencil-with-transversal
configuration is straightenable by choosing an affine chart with
$d$ at infinity.  Write its normal rows as $a_i=(w_i,z_i)$,
$a_j=(w_j,z_j)$, $a_k=(w_k,z_k)$ with $w_\ell\in\mathbb R^2$.
After reorienting $i,j$ by $\varepsilon,\eta$, their vectors
$\varepsilon w_i$ and $\eta w_j$ lie on the same ray. Write
$w_j=\lambda w_i$, where $\lambda\ne0$ and
$\operatorname{sgn}\lambda=\varepsilon\eta$. Elementary expansion gives
\[
 \det(a_i,a_j,a_k)=-(z_j-\lambda z_i)\det(w_i,w_k).
\]
On a large circle of radius $R$, the signed endpoints are near
$J(\varepsilon w_i)/\|w_i\|=J(\eta w_j)/\|w_j\|$, where $J$
rotates by $\pi/2$. Their angular offsets are ordered by
$\varepsilon z_i/\|w_i\|$ and $\eta z_j/\|w_j\|$.
Consequently the sign of the second offset minus the first is the sign of
$-\eta\chi_M(i,j,k)q(i,k)$, as in
\eqref{eq:polygon-expansion}.  The order is geometric and independent of
the choice of transversal $k$, so the formula is independent of $k$
and is transitive for the whole packet.  Only the local pencil and
transversal were straightened, which proves the formula for arbitrary
pseudoline arrangements.

For a small example, take old equations $x=0$, $x=z$, $y=0$
and the new line $z=0$. The three old normal vectors have determinant
$+1$. The parallel packet calculation gives the word
$(2,1,3,-1,-2,-3)$; its half-turn sends $2$ to $-1$.
Label negation and positional half-turn therefore differ as soon as
a nontrivial packet occurs.
\end{proof}

\subsection{The full image criterion and a local obstruction}
The following criterion combines the global chirotope restrictions
with the literal endpoint order. Both pieces are needed.

\begin{proposition}[The complete signed-polygon image]\label{prop:polygon-image}

A cyclic signed word $W$ is the deployed boundary polygon of a
ordinary relative position for the prescribed $\Sigma$ if and
only if the following finite tests hold, with the calibration just
specified:
\begin{enumerate}
 \item $\chi_{q_W}$ satisfies the unsigned basis axiom and
       \eqref{eq:polygon-gp};
 \item no row of $q_W$ is identically zero;
 \item $W$ agrees, up to a cyclic shift, with
       $\operatorname{Exp}_M(q_W)$.
\end{enumerate}
If $W$ is given without a cyclic orientation, perform these steps with
either orientation; the resulting extension is taken modulo
reorientation of $d$.  The tests include arbitrary vertex
multiplicities and any number of vertices on $d$.
\end{proposition}
\begin{proof}

For an actual relative position, cut along a narrow M\"obius band about
$d$.  Two remaining chords cross exactly when their intersections with
$d$ were distinct.  By the Jordan separation property for a proper arc
of a disk, they cross exactly when their endpoints alternate.  In that
case their signed cyclic order gives $\chi_N(i,j,d)$, by the
rank-two chirotope convention.  Thus the pair rule gives $q_W=q$.
The combined map is the chirotope of $N$, and it passes the two axioms.
A zero row at $i$ would say that $\{i,d\}$ has rank one: $d$ would be
parallel to $i$ in the oriented matroid, and hence the same projective
pseudoline.  This is excluded.  Finally the local disk computation in
Proposition~\ref{prop:packet-order} is exactly the expansion of the packets, so the third condition
holds.

Conversely, the first condition produces a rank-three oriented matroid
$N$ with chirotope $\chi_{q_W}$, whose deletion of $d$ is exactly $M$.
The second condition makes $d$ a nonloop and prevents it from being
parallel to any old element.  All old elements were already distinct
and nonparallel.  Therefore the projective topological representation
of $N$ consists of the prescribed number of distinct pseudolines.
Its deletion represents $M$, so the relative uniqueness/homeomorphism
comparison stated above carries that deletion to $\Sigma$ and carries
the new pseudoline to a relative position for $\Sigma$.  Proposition~\ref{prop:packet-order} computes
its boundary word as $\operatorname{Exp}_M(q_W)$; the third condition
identifies that word with $W$.  This proves sufficiency and exhausts
the image.  Reversal changes the contraction orientation, and no old
label is reoriented in this comparison.
\end{proof}

For comparison, consider the strongest natural local triangle test
on a stable candidate word. An adjacent pair
$(\varepsilon i,\eta j)$ is a candidate empty triangle if for every
$k\notin\{i,j\}$ the nonzero signs in
\begin{equation}\label{eq:polygon-triangle}
 \chi_M(i,j,k)q(i,j),\qquad
 \varepsilon q(i,k),\qquad \eta q(j,k)
\end{equation}
agree. For a valid simple arrangement these are the sides, relative
to the $k$th old chord, of the vertex and the two boundary endpoints.
Agreement says that this chord does not cut the triangle. Adjacency
excludes an endpoint in its base. This explains the local condition
without assuming that a candidate word is valid.

\begin{proposition}[Three local tests do not determine the image]\label{prop:polygon-obstruction}

The $2n$ weight rule, the cluster-reversal rule, and the presence of at
least three empty-triangle candidates, even tested against every old
line by \eqref{eq:polygon-triangle}, do not characterize the image.  This
fails for a simple arrangement of six straight lines.  The complete
replacement is Proposition~\ref{prop:polygon-image}.
\end{proposition}
\begin{proof}

Let $a_i=(1,t_i,t_i^2)$ for $t_i=i-1$, $1\le i\le6$, and let the old
projective line $i$ have equation $a_i\cdot x=0$.  For $i<j<k$,
\[
 \det(a_i,a_j,a_k)=(t_j-t_i)(t_k-t_i)(t_k-t_j)>0.
\]
Thus all old triples have sign $+1$, all old lines are distinct, and
no three old lines are concurrent.  In particular the arrangement is
non-pencil.  Consider the cyclic word
\[
 W=(1,2,-3,-4,-6,-5,-1,-2,3,4,6,5).
\]
It is stable and exactly antipodal: the second half is the negative of
the first. Its rank-two pair signs are the following upper triangle;
the other entries follow by alternation.
\[
\begin{array}{c|rrrrr}
i\backslash j&2&3&4&5&6\\ \hline
1&+&-&-&-&-\\
2&&-&-&-&-\\
3&&&+&+&+\\
4&&&&+&+\\
5&&&&&-
\end{array}
\]
The three adjacent arcs
\[
                (2,-3),\qquad(-1,-2),\qquad(3,4)
\]
pass \eqref{eq:polygon-triangle} for every other old line.  For the first
arc all four sign triples are $(-,-,-)$; for the other two arcs all
four sign triples are $(+,+,+)$.  The remaining nine adjacent arcs
fail, so there are exactly three candidates.  Give each marked arc
weight one.  In their cyclic order the three complementary arcs have
weights $4,1,4$, the numbers of boundary edges they contain.  The
total weight is $1+1+1+4+1+4=12=2n$. Every packet is a singleton,
so the cluster rule holds vacuously.

Nevertheless, in the combined map of \eqref{eq:polygon-combined}, take
$\lambda=5$, $a=1$, $b=3$, $c=6$, and $e=d$.  The three terms in
\eqref{eq:polygon-gp} are all $-1$.  For example they are
\[
 \chi_M(5,1,3)q(5,6)=-1,\quad
 -\chi_M(5,1,6)q(5,3)=-1,\quad
 q(5,1)\chi_M(5,3,6)=-1.
\]
This violates the chirotope axiom, so no pseudoline in the prescribed
arrangement has the candidate word.  The displayed sign computations check all four other old lines
at each candidate and exhibit the global obstruction directly.
Thus the counterexample concerns this prescribed arrangement.
\end{proof}

\subsection{Recovering weights and cyclic incidences}
There are two graphs to distinguish. The directed specialization
graph has all ordinary positions and the $n$ parallel positions as
vertices. Its arrows are ordinary one-vertex specializations and
arrows from one-zero positions to their adjacent parallel positions.
The dual graph of $\mathcal X_2$ has only ordinary vertices; besides
ordinary transitions, it has exceptional sweeps pairing the two
normal choices at each old edge and old vertex. We reconstruct both,
and then the endpoint transport for the surface. A one-zero position
will also be called semi-stable.

\begin{proposition}[Reconstruction from the full polygon image]\label{prop:polygon-reconstruction}

For a fixed labelled non-pencil arrangement $\Sigma$, its full image of
deployed polygons determines all ordinary relative positions,
all ordinary specialization arrows, their directions and weights, and
their cyclic incidences at each position.  Adjoining the known $n$
parallel positions recovers the corridor arrows and their
weights.  The same data determine all ordinary and exceptional sides of the
polygonal faces. With the prescribed endpoint transport, they determine
the weighted cellular map $\mathcal X_2$, up to label-preserving
isomorphism.
The reconstructed graph alone does not assert that every face has
order four; that is the separate local-link assertion of Section~\ref{sec:links}.
\end{proposition}
\begin{proof}

Proposition~\ref{prop:polygon-injective} reconstructs the extension at every polygon.  Its signed
cocircuit localization is obtained by setting
$\sigma(C_{ij})=q(i,j)$, where $C_{ij}(k)=\chi_M(i,j,k)$.
Equal cocircuit representatives give equal values because the
combined map is an extension chirotope.  A projective old vertex
$v$ is a pair $\{C_v,-C_v\}$ of cocircuits; write
$F_v=C_v^0\subset E$ for the set of old lines through it, and
$m_v=|F_v|$.  An ordinary immediate specialization is obtained by
replacing the two values at $\pm C_v$ by zero and leaving every other
localization value unchanged, provided the resulting localization
passes the extension criterion and is not parallel.  This rule
is finite and orientation-independent: simultaneous negation of all
values just reorients $d$.

We verify the geometric equivalence of the rule.  An elementary sweep
to $v$ moves no other old vertex across $d$, so it has exactly this
effect on the localization.  Conversely, if the rule yields an
extension $N'$, represent it with its new line through $v$.  Choose a
disk about $v$ containing no other old vertex.  Resolving the new line
to either side inside that disk gives the two localizations obtained
by replacing the zero at $C_v$ by the two possible signs; all other
coordinates are unchanged.  Exactly one has the given nonzero sign
of the original extension $N$.  Uniqueness of the extension therefore
identifies that local resolution with $N$.  This comparison is a
homeomorphism preserving the old arrangement cellwise, so the local
resolution proves that the required empty fan exists for the original
relative position.  In a cut-open disk it is precisely the
specialization triangle, subdivided by the $m_v$ concurrent old
lines.  There are $m_v-1$ boundary edges in its base, which proves
that the arrow weight is $m_v-1$.

At $N'$, the class $F_v$ is visible directly from the endpoint
nonalternation relation.  Its two boundary packets give the two
generization intervals.  Reversing one packet performs the
corresponding resolution; the other packet stays fixed.  At $N$, the
same marked interval is the specialization base.  Thus the signed
boundary word specifies the interval occupied by each incident
ordinary arrow, not only its two end vertices.  Their circular order
is the required local rotation, and the lengths of the complementary
intervals give all corner weights.  The intervals partition the
$2n$ boundary edges between transition bases and corner intervals,
which also proves the $2n$ weight identity.

For a parallel position $i$, use the parallel extension
$N_i$ with $d$ parallel to $i$.  Its localization is
$\sigma_i(C)=C(i)$, and its contraction entries are
$q_i(j,k)=\chi_M(j,k,i)$.  Let $N$ be semi-stable, with its unique old
vertex $v$ on line $i$.  There is a corridor arrow $N\to N_i$ exactly
when, after choosing the sign of the new element,
\[
          \sigma_i(C)\ne0 \quad\Longrightarrow\quad
          \sigma_N(C)=\sigma_i(C) .                     \label{eq:polygon-corridor}
\]
Necessity is immediate by collapsing the corridor.  For sufficiency,
contract $i$.  The cyclic rank-two localization criterion shows that,
with the zero at $v$ prescribed, there are just the two choices
corresponding to the two local sides of $i$ at $v$.  A narrow strip
about $i$, with the new line crossing $i$ at $v$ and following either
side elsewhere, realizes these choices.  Its localization away from
$i$ is $C(i)$; on $i$ it is the prescribed rank-two pattern.
Therefore \eqref{eq:polygon-corridor} and uniqueness identify $N$ with one
of these narrow-strip positions.  This is the required corridor.
The $n-m_v$ lines not through $v$ cross the corridor as disjoint
transversal chords, dividing its side into $n-m_v+1$ edges.  Hence
the ordinary corridor weight is $n-m_v+1$; its modified weight at
the parallel vertex is $m_v-1$, and their sum is $n$.
The order of the corridor flags at $i$ is the order of the old
vertices on the twofold cover of $i$, also determined by $M/i$.

To recover the surface graph, use the boundary circle of the parallel
polygon at $i$, identified with the rank-two extension circle of $M/i$.
Its alternating near positions $S_v^\pm,U_e^\pm$ and their cyclic
order are determined by that contraction. Pair the two normal choices
at every $v$ and $e$ to recover all exceptional sweeps, including
those at stable positions. Preserve the endpoint names prescribed in
Section~\ref{sec:construction}. Thus the parallel vertex of the
specialization graph is replaced by its full band data; it is not
itself a vertex of the dual surface graph.

Finally, a graph with cyclic orders at its vertices and specified
identifications of the two local orientation choices across its edges
is a signed rotation system.  Thicken each vertex to a disk and each
edge to a band with that identification.  Its boundary cycles are
determined by following the finite permutations of flags; attaching
one disk to each such cycle gives the unique cellular map with these
data.  Transition transport is part of the construction being
reconstructed, and its flag identifications are obtained from the
labelled local sweep or the corridor.  This establishes the asserted
reconstruction without deducing quadrangularity or embeddability of
the original large disks from the mere existence of the signed
rotation system.
\end{proof}


\subsection{Straight positions in one fixed realization}
A straight realization illustrates why the surface theorem uses all
pseudoline extensions. For the six vectors in the second arrangement
of Theorem~\ref{thm:examples}, let $v_{ij}=a_i\times a_j$.
A new straight line with coefficient vector $p$ passes through the
old vertex $v_{ij}$ exactly when $p\cdot v_{ij}=0$.
Thus straight positions are the cells of a fifteen-line arrangement
in the dual projective plane, with the six points $p=a_i$ excluded.

Exact intersections give 42 double vertices, one triple vertex
$p=(0,0,1)$, and six vertices of multiplicity five, namely the old
line vectors. Therefore the dual arrangement has
\[
 V=49,\qquad E=42\cdot2+3+6\cdot5=117,\qquad
 F=1-V+E=69.
\]
Here $E$ counts the distinct vertices on each projective dual line;
a projective circle with $r$ marked points has $r$ open edges.
The arrangement is a disk cellulation, so its projective Euler
characteristic gives the displayed value of $F$.
The ordinary straight positions have counts
\[
 (c_0^{\rm str},c_1^{\rm str},c_2^{\rm str},c_3^{\rm str})
       =(69,117,42,1),
\]
and total $229$. Their weighted zero count is
$117+2\cdot42+3=204$. These differ from the pseudoline counts
$(80,150,75,10)$ and total $315$ of Section~\ref{sec:six}.
The number of vertex--face sectors in the dual arrangement is
$2E=234$.
At the straight triple position, three concurrent dual walls have
six angular sectors, while the independent topological resolutions
of three old vertices give eight stable positions. Allowing the old
realization itself to move would answer a different question.

\begin{proposition}[The straight parameter surface]\label{prop:straight-blowup}
Let $B$ be the real blowup of the dual projective plane at the six
old-line points $a_1,\ldots,a_6$ of the displayed configuration.
It compactifies the space of ordinary straight lines and is a
closed surface homeomorphic to $\#^7\mathbb {RP}^2$.
Its Euler characteristic is $-5$. The strict transforms of the
fifteen dual lines, together with the six exceptional circles,
give a cellulation with
\[
 (f_0,f_1,f_2)=(73,147,69).
\]
There is also a cellulation whose closed faces are indexed by all
$229$ ordinary straight positions and whose counts are
\[
 (f_0,f_1,f_2)=(234,468,229).
\]
Every vertex of this latter cellulation has order four. It is obtained
by the explicit cubical subdivision and band replacement below.
The ordinary straight positions, considered as a subset of the
pseudoline extensions, are not closed under single-vertex resolutions.
In particular $B$ is not obtained by merely discarding the nonstraight
positions from the pseudoline reflection gluing.
\end{proposition}
\begin{proof}
The local real blowup replaces the origin in a disk by its unoriented
tangent directions. A closed neighborhood of its exceptional circle
has the explicit presentation
\[
 [-\varepsilon,\varepsilon]\times[0,\pi]\big/
       (t,0)\sim(-t,\pi),
\]
with blowdown $(t,\theta)\mapsto
(t\cos\theta,t\sin\theta)$. It is a M\"obius band, and its boundary
maps homeomorphically to the old disk boundary. Thus blowing up a
real surface point replaces a disk by a M\"obius band, the usual
connected-sum description~\cite[Corollary~F.3.2, p.~465]{Mangolte}.
The six distinct centers have disjoint disk neighborhoods. Starting
with $\mathbb {RP}^2$ therefore gives seven crosscaps. Away from the
exceptional circles the blowdown is a homeomorphism to the dual plane
with the six centers deleted, exactly the space of ordinary straight
lines. This proves the compactification statement and its topology.

Each old-line point lies on five distinct dual lines. On blowing it
up, these meet its exceptional circle at five distinct points. The
six original vertices are therefore replaced by thirty vertices,
so the vertex number is $49-6+30=73$. Each strict transform of a dual
line is still a circle with the same number of marked points, hence
still contributes its original edges. Each exceptional circle adds
five edges, giving $117+30=147$. The sixty local sectors at the six
centers are simply truncated along their exceptional circles;
the interiors of the original sixty-nine disk faces are unchanged.
Their closures are the corresponding truncated polygons, giving
the stated cellulation and the independent check
$73-147+69=-5$.

We next construct the second cellulation, including its incidences.
For any face of the original dual arrangement, mark its center and
the midpoints of its boundary edges. Join the center to each midpoint.
This cuts the face into one quadrilateral for each of its corners:
its vertices are the old corner, the two adjacent edge midpoints,
and the face center. The resulting cubical subdivision $Q$ has one
vertex for every cell of the dual arrangement, two edges for each
old edge joining its midpoint to its endpoints, and two more for
each old edge joining its midpoint to its adjacent face centers.
Its quadrilaterals are the old vertex--face sectors. Thus
\[
 (V(Q),E(Q),F(Q))=(49+117+69,\,4\cdot117,\,2\cdot117)
                   =(235,468,234).
\]
The old line arrangement is a regular disk cellulation: its spherical
chambers are convex and the projective antipodal identification pairs
distinct chambers. In particular the quadrilaterals just described
are embedded, and each vertex star used next is an embedded disk.

An old-line point $a_i$ has ten incident edge germs and ten sectors,
since it lies on five dual lines. Its closed star in $Q$ consists
of ten quadrilaterals. Its boundary has twenty vertices alternating
between old edge midpoints and face centers. Label them cyclically
$b_0,\ldots,b_{19}$ so that $b_j$ and $b_{j+10}$ are opposite
directions around $a_i$. Replace this disk by the following
quadrangulated M\"obius band, with exactly the same boundary. Add
the ten edges
\[
             b_jb_{j+10}\qquad(0\le j<10)
\]
and the ten quadrilaterals
\[
 (b_j,b_{j+1},b_{j+11},b_{j+10})\qquad(0\le j<10),
\]
where indices are taken modulo $20$. To verify the topology, take
a strip of ten squares, and identify its two ends with reversed
transverse coordinate. Its boundary runs successively through all
twenty displayed vertices, its transverse edges are exactly the ten
added edges, and its squares have the displayed boundary cycles.
This description also proves that all its closed cells embed.
Changing the starting label or reversing the cyclic labelling gives
the same incidence prescription up to relabelling.

The replacement removes the central vertex, replaces its ten radial
edges by ten band edges, and replaces ten quadrilaterals by ten
quadrilaterals. Perform it at the six old-line vertices. The interiors
of these stars are disjoint; their boundary edges and vertices, if
shared, remain fixed. Equivalently, first replace six disjoint small
disks inside the stars by M\"obius bands and then absorb the intervening
annular collars into the displayed bands. This identifies the resulting
surface with the real blowup $B$ already constructed, including at
points shared by star boundaries. The resulting quadrangulation $Q_B$
has counts
\[
                  (V(Q_B),E(Q_B),F(Q_B))=(229,468,234).
\]
Its vertices are exactly the original dual vertices other than the
six old-line points, the original edge midpoints, and the original
face centers. These are precisely all the ordinary straight positions.

Take the cellular dual of $Q_B$, placing one dual vertex in each
quadrilateral and joining across each edge. A small dual disk about
each vertex of $Q_B$ is an embedded closed face: the surface link is
a circle and each incident quadrilateral meets that vertex only once.
These dual disks cover $B$ and are indexed by the ordinary straight
positions. Duality gives counts $(234,468,229)$, and each dual vertex
has order four because its primal cell is a quadrilateral. This proves
the second cellulation and its direct Euler calculation
$234-468+229=-5$.

For the last assertion, the old vertices
$v_{12}=[0:1:0]$, $v_{34}=[1:0:0]$ and $v_{56}=[1:-1:0]$
are distinct and collinear. A straight line through any two is
necessarily $z=0$, so it also passes through the third.
In contrast, a pseudoline can be changed in a small disk about
$v_{56}$ to miss that vertex while remaining fixed near
$v_{12}$ and $v_{34}$. The same holds with any of the three
vertices chosen. These are the single-vertex resolutions already
used in the cube construction. None is straight in the fixed
realization. Deleting the nonstraight positions would remove the
partners of these resolution intervals, so the prescribed closed
reflection gluing does not restrict to the straight subset.
The real blowup above is an explicitly defined parameter
compactification, and no such restriction is needed to construct it.
\end{proof}

\subsection{Proper chords and boundary tethers}
The endpoint type is essential when passing between different arc
descriptions of a disk. A \emph{boundary tether} is a tame embedded
arc with one endpoint in the disk interior and one on its boundary,
and no other point on the boundary.

\begin{proposition}[The marked-chord correspondence]
\label{prop:tether}
A proper chord with a marked interior point is equivalent to an
unordered pair of boundary tethers with that common interior endpoint,
otherwise disjoint, and with distinct boundary endpoints. A single
boundary tether is not a proper chord. For a family of such pairs,
the one-crossing and disjointness-class conditions in
Proposition~\ref{prop:disk-criterion} remain separate hypotheses.
\end{proposition}
\begin{proof}
Parametrize a proper chord by $c:[0,1]\longrightarrow\Delta$ and
write the marked point as $c(t_0)$, where $0<t_0<1$. The restrictions
to $[0,t_0]$ and $[t_0,1]$ are exactly the two required tethers.
Conversely, two embedded intervals meeting just at their common
interior endpoint concatenate to an embedded interval. Their distinct
remaining endpoints lie on the boundary, and all other points lie
inside the disk, so the union is a proper chord. These operations
are inverse up to a homeomorphism preserving the marked point and
the unordered pair; ordering the pair orients the chord.

For a concrete distinction, in $\Delta=[-1,1]^2$ the arc
$\gamma(t)=(t,0)$, $0\leq t\leq1$, is a tether, whereas
$c(t)=(2t-1,0)$ is a proper chord. The first has one endpoint on
$\partial\Delta$ and the second has two. A homeomorphism of pairs
preserves that number. Moreover the tether can be completed by an
arc to $(-1,0)$ or by an arc to $(0,1)$, so a single tether does not
determine its missing half. Applying the construction to several
pairs imposes no restriction on how their unions intersect; the
additional family hypotheses must therefore be checked separately.
\end{proof}
